\documentclass[../DoAn.tex]{subfiles}
\begin{document}

\begin{center}
    \Large{\textbf{ABSTRACT}}\\
\end{center}
\vspace{1cm}
This thesis presents ThePawsome, a pet-focused e-commerce system that combines online shopping workflows with a Vietnamese AI assistant. The problem comes from the fact that pet owners often have to choose food, accessories, and care information from fragmented sources, while each pet has different species, age, weight, health notes, and allergies. Conventional shopping websites can display products and process orders, but they provide limited personalized advice, weak reuse of trusted community knowledge, and insufficient protection against duplicated orders or overselling when many customers buy the same product variant.

The proposed solution follows a separated frontend/backend web architecture. The backend is implemented with FastAPI, asynchronous SQLAlchemy, PostgreSQL with pgvector, and Redis. The frontend uses Next.js App Router, TypeScript, Tailwind CSS, TanStack Query, and Zustand. The Catbot assistant is built with LangGraph tool-calling. It retrieves real products, curated knowledge, and eligible forum discussions before generating grounded responses, while applying AI safety checks and human handoff for sensitive cases. The e-commerce workflow supports variant-based catalog management, cart, guest checkout, promotions, post-purchase reviews, returns, SePay VietQR payment, temporary inventory reservation, and idempotency handling to reduce duplicate order creation. The final result is a maintainable monorepo covering customer, admin, forum, AI/RAG, and core backend test suites, providing a practical foundation for future deployment with real operational data.

\end{document}
